\section{Two-step extensions in higher dimensions}\label{sec:higher-dim}

Let $(X,\omega)$ be a compact connected K\"ahler manifold of complex
dimension $n\ge1$, with the conventions of Section~\ref{sec:intro}. Let
$F_1,\ldots,F_m$ be pairwise nonisomorphic stable vector bundles of common
slope, with Hermitian--Einstein metrics $h_i$. Partition the indices as
$I\sqcup J$, and let
\[
 0\longrightarrow S=\bigoplus_{i\in I}F_i\longrightarrow E\longrightarrow
 U=\bigoplus_{j\in J}F_j\longrightarrow0
\]
be an extension, with harmonic representative $N=\sum\beta_{ij}$. Here each
entry $\beta_{ij}$ is a harmonic form in $\Omega^{0,1}(\Hom(F_j,F_i))$.
Assume that the graph $\Gamma$ of nonzero entries is connected.

\begin{theorem}\label{thm:higher-dim}
Let $E$ be as above, with masses $M_i=V\rk F_i$ and conductances
$k_{ij}=\lVert(gNg^{-1})_{ij}\rVert_{L^2}^2$. Along the normalized flow from
any smooth initial metric on $E$:
\begin{enumerate}
\item every small eigenvalue satisfies the two-sided comparison
 \eqref{eq:transfer} with the diagonal flow, with $\delta_T\to0$, and the
 full small spectral projection converges to $\mathcal K$ in every
 $L^2\to C^k$ norm;
\item if $(\Gamma,(r_i))$ is nondegenerate, conclusions (1)--(4) of
 Theorem~\ref{thm:complete-spectrum} hold;
\item for two factors, $\tau\lambda_1(h(\tau))\to\frac12$, as in
 Theorem~\ref{thm:flow-two-factor}.
\end{enumerate}
Corollaries~\ref{cor:coefficient-transfer} and~\ref{cor:balanced} and the
statements of Section~\ref{sec:flow-examples} hold on $X$ whenever the
required bundles and nonzero harmonic entries exist there.
\end{theorem}

The fixed-metric results of Sections~\ref{sec:two-factor}
and~\ref{sec:several}, including Theorem~\ref{thm:orders}, also hold on $X$
for two-step extensions. Section~\ref{sec:two-factor} is already stated
there. For more than two socle layers the families need not be integrable
when $n\ge2$. On
$\Sigma_1\times\Sigma_2$, for instance, entries $\beta_{12}$ and $\beta_{23}$ pulled
back from the two factors have a nonzero product class in
$H^{0,2}(\Hom(F_3,F_1))$.

\begin{proof}
We check each input used on curves.

\emph{Integrability.} Harmonic forms are $\db_0$-closed, and
$\beta_{ij}\wedge\beta_{kl}$ requires $j=k$, which cannot occur since
$j\in J$ and $k\in I$. So $\db_0N+N\wedge N=0$, and the same holds for
$gNg^{-1}$ and every entrywise rescaling.

\emph{Mean curvature.} The Chern connection of $(\db_0+N,h_0)$ is
$A_0+N-N^*$. Since $\db_0^*N=0$, the K\"ahler identity
$[\Lambda,\partial_0]=\sqrt{-1}\,\db_0^*$ gives $\Lambda\partial_0N=0$, and
likewise for $\db_0N^*$. Hence
$K-\kappa\id=\sqrt{-1}\Lambda(N-N^*)\wedge(N-N^*)$. In a local unitary
coframe with $N=\sum_\nu n_\nu\bar\theta^\nu$, this is
$\sum_\nu[n_\nu,n_\nu^*]$, as in \eqref{eq:two-C}. Its block traces give the
same $p_i$ as on a curve.

\emph{Fixed-metric comparison.} The proofs of Theorem~\ref{thm:relative}
and Proposition~\ref{prop:subspace} use the cancellation
$D_0^*[N,x]=[\db_0^*N,x]=0$, the spectral gap of $\Delta_0$ on $QH$,
elliptic estimates, and Sobolev embedding. All of these hold on $X$.

\emph{Order comparison.} Propositions~4.3 and~4.5 of \cite{HKKP2018},
used in \eqref{eq:flow-barriers}, are stated and proved on compact
K\"ahler manifolds of any dimension, by a pointwise maximum principle.

\emph{Correction.} The first variation of mean curvature at a
Hermitian--Einstein metric is $\Delta_\partial=\Delta_{\db}$ on
endomorphisms: the difference is $[K,\cdot]$, which vanishes there. The
mixed terms in Lemma~\ref{lem:general-correction} are
$O(\mathfrak e^{3/2})$ uniformly. So that lemma and Lemma~\ref{lem:layer-energy} hold
unchanged.

\emph{Convergence of the flow.} Every Jordan--H\"older filtration of $E$
has graded object isomorphic to $S\oplus U$. A saturated stable subsheaf
of slope $\mu$ in $E$ is reflexive, hence isomorphic to some $F_k$, and
$\Hom(F_k,E)=0$ for $k\in J$ since each extension class is nonzero. So the
Harder--Narasimhan--Seshadri graded object is locally free, and its singular
set is empty. By Sibley--Wentworth \cite[Theorem~1.1]{SibleyWentworth2015},
the flow converges modulo unitary gauge, smoothly on all of $X$, to an
admissible Yang--Mills connection $A_\infty$ on a bundle isomorphic to
$S\oplus U$. They credit the smooth convergence to Hong and Tian
\cite{HongTian2004}, and the
identification of the limit to Jacob and Sibley
\cite{Jacob2015,Sibley2015}.

$A_\infty$ is smooth and integrable. Its Hermitian--Einstein summands have
constants given by their slopes, and every holomorphic summand of
$\bigoplus F_i$ has slope $\mu$. So $A_\infty$ is Hermitian--Einstein, and it
is gauge equivalent to the product connection of $h_0$. For each $k$, the
infimum over unitary gauges of the $C^k$ distance to $A_\infty$ tends to
zero. Choosing near-optimal gauges with $k$ tending slowly to infinity gives,
for every large $\tau$, holomorphic isometries from $(E,h(\tau))$ to
$(E_0,h_0,\db^{(\tau)})$ with $\db^{(\tau)}\to\db_0$ in $C^\infty$. This is
the input used on curves.

\emph{Matching.} In Lemmas~\ref{lem:flow-matching}
and~\ref{lem:socle-matching}, the transported extension entry $\alpha_T$ is
$\db$-closed: in the block computation of $D_T^2=0$, the product
$\alpha_T\wedge\alpha_T$ lands in $\Hom(U,S)\circ\Hom(U,S)=0$. Hodge theory
on $X$ gives $\alpha_T=H\alpha_T+\db v_T$ with $v_T\to0$ in $C^\infty$.
Uniqueness up to diagonal maps uses only $L^2$-orthogonality of harmonic
and exact forms and $H^0(\Hom(U,S))=0$.

The remaining arguments are finite-dimensional or use only these inputs.
\end{proof}

\begin{example}\label{ex:products}
Let $X=\Sigma\times Y$, where $\Sigma$ is a curve of genus at least two and $Y$ is a
connected compact K\"ahler manifold, with a product K\"ahler form.
\begin{itemize}
\item For stable $F$ on $\Sigma$, the pullback of a Hermitian--Einstein metric is
 Hermitian--Einstein, since $\Lambda_\omega$ of a form pulled back from $\Sigma$
 is $\Lambda_{\omega_\Sigma}$ of that form. By the K\"unneth formula
 $H^0(X,\End\pi^*F)=\CC$, so $\pi^*F$ is stable.
\item Nonisomorphic stable bundles of equal slope pull back to such
 bundles, and
 \[
  H^1(X,\Hom(\pi^*G,\pi^*F))\supset H^1(\Sigma,\Hom(G,F))\ne0.
 \]
\item Harmonic forms pull back to harmonic forms.
\end{itemize}
Every example of Section~\ref{sec:flow-examples} therefore occurs on $X$,
with arbitrary initial metrics on $X$. The same holds for the sharpness
families of Section~\ref{sec:slope}. For the explicit constants,
$V=1$ means $\operatorname{Vol}(\Sigma)\operatorname{Vol}(Y)=1$.
\end{example}

For extensions with more than two steps in dimension $n\ge2$, the
Maurer--Cartan equation prevents all entries from being harmonic. A natural
replacement is the gauge $\db_0^*N=0$. In it, the entry $N_{13}$ acquires
the cup-product correction $-\db_0^*\mathbb G(N_{12}\wedge N_{23})$, with
$\mathbb G$ the Green operator, which exists
only when $[N_{12}]\cup[N_{23}]=0$. We do not treat this case.
